% problem_id: S047-b21-Q2
% source_id: S047/S048 (numdam/aif3604.pdf)
% source_file: aif3604.pdf
% statement_locator: printed p. PDF p. 15 = printed p. 2518 (statement starts on that page, line "THEOREM G. — Suppose that p ∈ Sing_1(X) is ...", and the first two lines of the proof, "Proof. — Let I be the ideal ..." / "Suppose that ψ : U′ → U is toroidal ...", are on the same page)
% proof_locator: printed p. PDF pp. 15-18 = printed pp. 2518-2521 (proof runs from PDF p. 15 through PDF p. 18, ending in the end-of-proof square on printed p. 2521)
% representation: PDF; transcribed from rendered page images
% collected_by: carrier rule
% review_status: H2 by 2/2 independent reviewers (the miner claimed H3); 2/2 忠实
% status: complete (statement + author proof verbatim + reason + locators)
%
% =====================================================================================
%  human-proof-corpus candidate transcription  —  candidate Q2
%
%  Source paper : Steven Dale CUTKOSKY, "Corrigendum to 'A simpler proof of toroidalization
%                 of morphisms from 3-folds to surfaces'", Ann. Inst. Fourier, Grenoble
%                 74, 6 (2024) 2505-2521.  DOI 10.5802/aif.3604
%                 file: /disks/sata1/yupeng/human-proof-corpus/source-cache/registry-sources/
%                       numdam/aif3604.pdf   (18 PDF pages, 1411602 bytes)
%                 Page map: printed page = 2503 + PDF page  (PDF 1 = journal cover, no
%                 printed number; printed 2505 = PDF 2).  All 18 printed pages 2505-2521.
%
%  COLLECTED UNIT (the object of this transcription, and the only new complete
%  theorem-plus-proof unit of this paper that had not yet been collected):
%                 THEOREM G  —  statement printed p. 2518 (= PDF p. 15), complete proof
%                 printed pp. 2518-2521 (= PDF pp. 15-18), ending in the end-of-proof
%                 square on printed p. 2521.
%
%  NOT part of the collected unit (clearly delimited CONTEXT blocks, transcribed only
%  because the unit cannot be read without them):
%                 [CONTEXT A]  form (A) of the revised Definition 3.3   (printed p. 2507)
%                 [CONTEXT B]  the construction of omega'                 (printed p. 2518)
%                 [CONTEXT C]  the three errata line-citations that send
%                              the reader of [3] to Theorem G            (printed p. 2521)
%
%  Method (per the project discipline): pdftotext -layout was used ONLY to locate text;
%  every transcribed character below was read from page images rendered with
%  pdftoppm -r 150/300/600/1200 -png and inspected with read_image.  Verbatim: no
%  rewriting, no completion, no silent correction.  As-printed oddities are transcribed
%  as printed and flagged with  %% [NOTE: ...]  at the exact spot; they are collected
%  again in the closing "Transcription notes" section.  No [ILLEGIBLE: ...] marker was
%  needed anywhere: nothing on the transcribed pages is illegible.
%
%  Cross-page markers:  %% --- page <printed page> (PDF p. <n>) ---
% =====================================================================================
\documentclass[11pt]{article}
\usepackage[T1]{fontenc}
\usepackage[utf8]{inputenc}
\usepackage{amsmath,amssymb,amsfonts,amsthm}
\usepackage[margin=1in]{geometry}

\newcommand{\fk}{\mathfrak{k}}                 % the ground field, printed as fraktur k
\newcommand{\OO}{\mathcal{O}}                  % local ring (as printed: no hat)
\newcommand{\OH}{\widehat{\mathcal{O}}}        % completion (as printed: with hat)
\newcommand{\Dbar}{\overline{D}}               % D with a bar
\newcommand{\ztil}{\widetilde{z}}              % z with a tilde
\newcommand{\tbar}{\bar{t}}                    % t with a bar
\newcommand{\Sing}{\operatorname{Sing}}
\newcommand{\ord}{\operatorname{ord}}

\setlength{\parindent}{0pt}
\setlength{\parskip}{0.5em}

\begin{document}

\begin{center}
{\large\bfseries Candidate Q2 --- transcription}\\[0.3em]
{\scshape Cutkosky, Corrigendum to ``A simpler proof of toroidalization of morphisms\\
from 3-folds to surfaces'', Ann.\ Inst.\ Fourier 74, 6 (2024) 2505--2521}\\[0.3em]
{\small collected unit: \textsc{Theorem G} (printed pp.\ 2518--2521 = PDF pp.\ 15--18)}
\end{center}

\hrule
\vspace{1em}

% =====================================================================================
\section*{0.\ \ What is transcribed, and in what order}
% =====================================================================================

\begin{itemize}\itemsep0.2em
\item \textbf{[CONTEXT A]} printed p.\ 2507 (= PDF p.\ 4): item (4) of the revised
      Definition 3.3, i.e.\ the definition of the form (A) to which Theorem G applies.
\item \textbf{[CONTEXT B]} printed p.\ 2518 (= PDF p.\ 15): the paragraph in which the
      errata constructs the function $\omega'$ used both in (A) and in Theorem G.
\item \textbf{[THE COLLECTED UNIT]} printed pp.\ 2518--2521 (= PDF pp.\ 15--18):
      \textsc{Theorem G}, statement and complete proof, verbatim.
\item \textbf{[CONTEXT C]} printed p.\ 2521 (= PDF p.\ 18): the three remaining errata
      items, which are the places where the corrected text of [3] must now cite
      ``Theorem 4.1, 4.2 or G'' --- evidence of Theorem G's role, not part of the unit.
\item \textbf{Transcription notes}: as-printed oddities, reading record, and what was
      deliberately not transcribed.
\end{itemize}

% =====================================================================================
\section*{[CONTEXT A] \ \ form (A) of the revised Definition 3.3 \quad (printed p.\ 2507 = PDF p.\ 4)}
% =====================================================================================

\emph{The errata announces this block on printed p.\ 2506 with:} ``Page 877: Definition
3.3 should be modified as follows.'' \emph{Transcribed from the page image:}

\medskip
(4) $p$ is a 1-point, and we have an expression (2.1) with

(A)\quad $F = \tau_0 z^m + \tau_2 x^{r_2} y z^{m-2} + \cdots + \tau_{m-1} x^{r_{m-1}} y z + \tau_m x^{r_m} y + x^t \Omega$

\smallskip
where $\tau_0 \in \OH_{X,p}$ is a unit, $\tau_i \in \OH_{X,p}$ are units (or zero) for
$2 \leqslant i \leqslant m$, $\Omega \in \OH_{X,p}$, $\tau_i \neq 0$ for some $i$ with
$2 \leqslant i \leqslant m$ and $t > \omega'(m, r_2, \dots, r_m)$ (where we set $r_i = 0$
if $\tau_i = 0$). Further, $r_i > 0$ if $\tau_i \neq 0$.

\medskip
$X$ is 3-prepared if $X$ is 3-prepared for all $p \in X$.

% =====================================================================================
\section*{[CONTEXT B] \ \ the construction of $\omega'$ \quad (printed p.\ 2518 = PDF p.\ 15)}
% =====================================================================================

\emph{The line preceding it, and the whole paragraph, transcribed from the page image. It
stands between the end of the errata's newly written detailed proof of Theorem 4.2
(end-of-proof square, printed p.\ 2518) and the statement of Theorem G:}

\medskip
Before the statement of Theorem 4.3 on page 898, add the following:

We construct the function $\omega'(m, r_2, \dots, r_m)$ in a similar way to the
construction of $\omega$. Let $I$ be the ideal in $\fk[x, z]$ generated by $z^m$ and
$x^{r_i} z^{m-i}$ for all $i$ such that $2 \leqslant i \leqslant m$ and $\tau_i \neq 0$.
We define $\omega'(m, r_2, \dots, r_m)$ as we define $\omega$, except we allow $i$ to
range within $2 \leqslant i \leqslant m$.

% =====================================================================================
\section*{[THE COLLECTED UNIT] \ \ \textsc{Theorem G} and its complete proof}
% =====================================================================================

%% --- page 2518 (PDF p. 15) ---

\medskip
\textsc{Theorem G.} --- \emph{Suppose that $p \in \Sing_1(X)$ is a 1-point and $X$ is
3-prepared at $p$. Let $x, y, z$ be permissible parameters at $p$ giving a form (A) at
$p$. Let $U$ be an \'etale cover of an affine neighborhood of $p$ in which $x, y, z$ are
uniformizing parameters. Then $xz = 0$ gives a toroidal structure $\Dbar$ on $U$.}

\emph{There is (after possibly replacing $U$ with a smaller neighborhood of $p$) a
unique, minimal toroidal morphism $\psi : U' \to U$ with respect to $\Dbar$ which has the
property that $U'$ is nonsingular, 2-prepared and $\Gamma_D(U') < \sigma_D(p)$. This map
$\psi$ factors as a sequence of permissible blowups $\pi_i : U_i \to U_{i-1}$ of sections
$C_i$ over the two curve $C$ of $\Dbar$. $U_i$ is 1-prepared for $U_i \to S$. We have
that the curve $C_i$ blown up in $U_{i+1} \to U_i$ is in $\Sing_{\sigma_D(p)}(U_i)$ if
$C_i$ is not a 2-curve of $D_{U_i}$, and that $C_i$ is in $\Sing_1(U_i)$ if $C_i$ is a
2-curve of $D_{U_i}$.}
%% [NOTE: as printed: "of sections $C_i$ over the two curve $C$ of $\Dbar$" — singular
%%  "the two curve"; elsewhere in the paper this object is written "2-curve". Kept verbatim.]

\medskip
\emph{Proof.} --- Let $I$ be the ideal in $\Gamma(U, \OO_X)$ generated by
\[
z^{r_m} \ \text{and}\ \{ x^{r_i} z^{m-i} \mid 2 \leqslant i \leqslant m \ \text{and}\ \tau_i \neq 0 \}.
\]
Suppose that $\psi : U' \to U$ is toroidal for $\Dbar$ and $U'$ is nonsingular. Let
$\Dbar' = \psi^{-1}(\Dbar)$.

%% --- page 2519 (PDF p. 16) ---

The set of 2-curves of $\Dbar'$ is the disjoint union of the 2-curves of $D_{U'}$ and the
2-curve which is the intersection of the strict transform of the surface $z = 0$ on $U'$
with $D_{U'}$. $\psi$ factors as a sequence of blow ups of 2-curves of (the preimage of)
$\Dbar$. We will verify the following three statements, from which the conclusions of the
theorem follow.

\begin{list}{}{\setlength{\leftmargin}{4.5em}\setlength{\itemindent}{-4.5em}}
\item[(c11)] If $q \in \psi^{-1}(p)$ and $I\OO_{U',q}$ is principal, then
      $\sigma_D(q) < \sigma_D(p)$.
\item[] \hspace*{4.5em} In particular, $\sigma_D(q) < \sigma_D(p)$ if $q$ is a 1-point of
      $\Dbar'$.
\item[(c12)] If $C'$ is a 2-curve of $D_{U'}$, then $U'$ is prepared at
      $q = C' \cap \psi^{-1}(p)$\\
      \hspace*{4.5em} if and only if $\sigma_D(q) < \infty$\\
      \hspace*{4.5em} if and only if $I\OO_{U',q}$ is principal\\
      \hspace*{4.5em} if and only if $U'$ is prepared at all $q' \in C'$ in a
      neighborhood of $q$.
\item[(c13)] If $C'$ is the 2-curve of $\Dbar'$ which is the intersection of $D_{U'}$
      with the strict transform of $\ztil = 0$ in $U'$, then $\sigma_D(q) \leqslant
      \sigma_D(p)$ if $q = C' \cap \psi^{-1}(p)$, and $\sigma_D(q') = \sigma_D(q)$ for
      $q' \in C'$ in a neighborhood of $q$.
\end{list}
%% [NOTE: (c11) and (c12) print the ideal as $I\OO_{U',q}$ (script O, no hat), whereas the
%%  working passages of the proof below print $I\OH_{U',q}$ (with hat). Both spellings are
%%  kept as printed.]
%% [NOTE: (c13) prints the surface as $\ztil = 0$ (z with tilde), whereas the corresponding
%%  "final case" three paragraphs below prints $z = 0$ (no tilde). Both kept as printed.]

Suppose that $q \in \psi^{-1}(p)$ is a 1-point for $\Dbar'$. Then $I\OH_{U',q}$ is
principal. At $q$, we have permissible parameters $x_1, y, z_1$ defined by
\begin{equation}
\tag{c14}
x = x_1^{a_1}, \qquad z = x_1^{b_1} (z_1 + \alpha)
\end{equation}
for some $a_1, b_1 \in \mathbb{Z}_+$ and $0 \neq \alpha \in \fk$. Substituting into (A),
we have
\[
u = x_1^{a a_1}, \qquad v = P(x_1^{a_1}) + x_1^{b a_1} G
\]
where
\begin{align*}
G = {} & \tau_0 x_1^{b_1 m} (z_1 + \alpha)^m + \tau_2 x_1^{a_1 r_2 + b_1 (m-2)} y (z_1 + \alpha)^{m-2} + \cdots \\
       & + \tau_{m-1} x_1^{a_1 r_{m-1} + b_1} y (z_1 + \alpha) + \tau_m x_1^{a_1 r_m} y + x_1^{a_1 t} \Omega .
\end{align*}

Let $x_1^s$ be a local generator of $I\OH_{U',q}$. We have that $a_1 t > s$ by our construction
of $\omega'(m, r_2, \dots, r_m)$. Let $G' = \frac{G}{x_1^s}$.
%% [NOTE: the source breaks this word across two lines as "con-" / "struction"; the
%%  line-break hyphen is not part of the text and is not reproduced.]

If $z^m$ is a local generator of $I\OH_{U',q}$, then $G'$ has an expansion
\[
G' = \tau' (z_1 + \alpha)^m + x_1 \Omega_1 + y \Omega_2
\]
where $0 \neq \tau' = \tau(0, 0, 0) \in \fk$, $g_2, \dots, g_m \in \fk$ and
$\Omega_1, \Omega_2 \in \OH_{U',q}$. We have $\ord(G'(0, 0, z_1)) \leqslant m - 1$.
Setting $F' = G' - G'(x_1, 0, 0)$ and $P'(x_1) = P(x_1^{a_1}) + x_1^{b a_1 + b_1 m}
G'(x_1, 0, 0)$, we have an expression
\[
u = x_1^{a a_1}, \qquad v = P'(x_1) + x_1^{b a_1 + b_1 m} F'
\]
of the form of (2.1). Thus $U'$ is 2-prepared at $q$ with $\sigma_{D'}(q) < m - 1 =
\sigma_D(p)$.
%% [NOTE: as printed, the $g_2, \dots, g_m \in \fk$ in this sentence are not used by the
%%  expansion displayed immediately above it (that expansion has no $g$'s); the phrase is
%%  carried over from the parallel passage of the proof of Theorem 4.2 (printed p. 2516),
%%  where the expansion does contain $g_2(z_1+\alpha)^{m-2} + \cdots + g_{m-1}(z_1+\alpha)$.
%%  Kept verbatim; no correction attempted.]

%% --- page 2520 (PDF p. 17) ---

Suppose that $z^m$ is not a local generator of $I\OH_{U',q}$, but there exists some $i$
with $2 \leqslant i \leqslant m$ such that $x^{r_i} z^{m-i}$ is a local generator of
$I\OH_{U',q}$. Let $h$ be the smallest $i$ with this property. Then $G'$ has an
expression
\[
G' = g_h (z_1 + \alpha)^{m-h} y + \cdots + g_m y + x_1 \Omega_1 + y_1^2 \Omega_2
\]
for some $g_i \in \fk$ with $g_h \neq 0$ and $\Omega_1, \Omega_2 \in \OH_{U',q}$. We have
that $U'$ is 2-prepared at $q$ with $\sigma_D(q) \leqslant (m - h + 1) - 1 < m - 1 =
\sigma_D(p)$ since $h \geqslant 2$.
%% [NOTE: as printed, the last term is $y_1^2 \Omega_2$ — with a subscript 1 on $y$ —
%%  although no $y_1$ is introduced in this (1-point) case: the coordinates here are
%%  $x_1, y, z_1$. The parallel passage of the proof of Theorem 4.2 (printed p. 2516)
%%  prints $y\Omega_2$ there. Kept verbatim.]

Now suppose that $q \in \psi^{-1}(p)$ is a 2-point for $D_{U'}$. We have permissible
parameters $x_1, y, z_1$ in $\OH_{U',q}$ such that
\begin{equation}
\tag{c15}
x = x_1^{a_1} z_1^{b_1}, \qquad z = x_1^{c_1} z_1^{d_1}
\end{equation}
with $a_1, b_1 > 0$ and $a_1 d_1 - b_1 c_1 = \pm 1$. Substituting into (A), we have
\[
u = x_1^{a_1 a} z_1^{b_1 a}, \qquad v = P(x_1^{a_1} z_1^{b_1}) + x_1^{a_1 b} z_1^{b_1 b} G
\]
where
\begin{align*}
G = {} & \tau_0 x_1^{c_1 m} z_1^{d_1 m} + \tau_2 x_1^{r_2 a_1 + c_1 (m-2)} z_1^{r_2 b_1 + d_1 (m-2)} y + \cdots \\
       & + \tau_{m-1} x_1^{r_{m-1} a_1 + c_1} z_1^{r_{m-1} b_1 + d_1} y + \tau_m x_1^{a_1 r_m} z_1^{b_1 r_m} y + x_1^{a_1 t} z_1^{b_1 t} \Omega .
\end{align*}

Let $C'$ be the 2-curve of $D_{U'}$ containing $q$. The three statements
$\sigma_D(q) < \infty$, $\sigma_D(q) = 0$ and $I\OO_{U',q}$ is principal are equivalent.
Further, we have that $\sigma_D(q') = \sigma_D(q)$ for $q' \in C'$ in a neighborhood of
$q$.

Suppose that $I\OO_{U',q}$ is principal and let $x_1^s z_1^{\tbar}$ be a local generator of
$I\OH_{U',q}$. We have that $a_1 t > s$ and $b_1 t > \tbar$ by our construction of
$\omega'(m, r_2, \dots, r_m)$. Let $G' = G / x_1^s z_1^{\tbar}$. We have that
\[
u = (x_1^{a_1} z_1^{b_1})^a, \qquad v = P(x_1^{a_1} z_1^{b_1}) + x_1^{a_1 b + s} z_1^{b b_1 + \tbar} G'
\]
has the form (2.2), since we have made a monomial substitution in $x$ and $z$. Since
$z^m$ or $x^{r_i} z^{m-i}$ for some $i$ is a local generator of $I\OH_{U',q}$, we have
that $\ord G'(0, y, 0) \leqslant 1$. We thus have that $U'$ is prepared at $q$.

The final case is when $q \in \psi^{-1}(p)$ is on the 2-curve $C'$ of $\Dbar'$ which is
the intersection of $D_{U'}$ with the strict transform of $z = 0$ in $U'$. Then there
exist permissible parameters $x_1, y, z_1$ at $q$ such that
\begin{equation}
\tag{c16}
x = x_1, \qquad z = x_1^{b_1} z_1
\end{equation}
for some $b_1 \in \mathbb{Z}_+$. The equations $x_1 = z_1 = 0$ are local equations of
$C'$ at $q$. Let
\[
s = \min \{ b_1 m,\ r_i + b_1 (m - i) \ \text{with}\ \tau_i \neq 0 \ \text{for}\ 2 \leqslant i \leqslant m \}.
\]

%% --- page 2521 (PDF p. 18) ---

We have that $t > s$ by our construction of $\omega'(m, r_2, \dots, r_m)$. We have an
expression of the form (2.1) at $q$,
\begin{align*}
u &= x_1^a \\
v &= P(x_1^a) + x_1^{a b + s} G'
\end{align*}
with
\begin{align*}
G' = {} & \tau_0 x_1^{b_1 m - s} z_1^m + \tau_2 x_1^{r_2 + b_1 (m-2) - s} z_1^{m-2} y + \cdots \\
        & + \tau_{m-1} x_1^{r_{m-1} + b_1 - s} z_1 y + \tau_m x_1^{r_m} y + x_1^{t - s} \Omega .
\end{align*}
We see that $\sigma_D(q) \leqslant \sigma_D(p)$ (with $\sigma_D(q) < \sigma_D(p)$ if
$s = r_i + b_1 (m - i)$ for some $i$ with $2 \leqslant i \leqslant m$) and
$\sigma_D(q') = \sigma_D(q)$ for $q'$ in a neighborhood of $q$ on $C'$.

Suppose that $I\OO_{U',q}$ is principal. Let $h$ be the largest $i$ such that $\tau_i \neq 0$
in (A). Then $x^{r_h} z^{m-h}$ is the local generator of $I\OH_{U',q}$. Since
$2 \leqslant h \leqslant m - 1$, we have that
\[
\sigma_D(q) = (m - h + 1) - 1 < m - 1 = \sigma_D(p). \qquad \square
\]
%% [NOTE: as printed, "Since $2 \leqslant h \leqslant m-1$". Here $h$ has just been defined as
%%  the LARGEST $i$ with $\tau_i \neq 0$ in (A), and (A) allows $i$ up to $m$, so $h = m$ is
%%  not excluded by that definition; the bound "$h \leqslant m-1$" is the one that holds in
%%  the parallel sentence on printed p. 2518 (proof of Theorem 4.2), where $h$ ranges over
%%  $2 \leqslant i \leqslant m-1$ because (3.7) has no $i = m$ term. With $h = m$ the
%%  displayed inequality still holds, since $\sigma_D(q) = (m-m+1)-1 = 0 < m-1$ for
%%  $m > 1$. Kept verbatim; recorded, not corrected. (The companion text of [3] at
%%  printed p. 896, which the errata also repairs, is exactly about splitting the cases
%%  $\tau_m \neq 0$ / $\tau_m = 0$ for this same generator claim.)]

% =====================================================================================
\section*{[CONTEXT C] \ \ the three errata items that cite Theorem G \quad (printed p.\ 2521 = PDF p.\ 18)}
% =====================================================================================

\emph{These three lines stand immediately after the end-of-proof square of Theorem G.
Transcribed from the page image; they are the places in [3] at which the corrected text
must send the reader to the new theorem:}

\medskip
Page 907, lines 19--20 and line 29: should be ``(3.6), (3.7) or (A)''

Page 908, line 9: should be ``Theorem 4.1, 4.2 or G''.

Page 918, lines 16 and $-2$: should be ``Theorem 4.2 and G''

\vspace{1em}
\hrule
\vspace{1em}

% =====================================================================================
\section*{Transcription notes}
% =====================================================================================

\subsection*{1.\ As-printed oddities, transcribed verbatim and flagged (never silently corrected)}
\begin{enumerate}\itemsep0.25em
\item \textbf{p.\ 2518, Theorem G, line 5 of the statement:} ``of sections $C_i$ over the
      two curve $C$ of $\Dbar$'' --- singular ``the two curve'' (the paper's usual word is
      ``2-curve''; cf.\ the two preceding occurrences of ``2-curve'' on p.\ 2519).
\item \textbf{p.\ 2519, in the proof:} ``where $0 \neq \tau' = \tau(0,0,0) \in \fk$,
      $g_2, \dots, g_m \in \fk$ and $\Omega_1, \Omega_2 \in \OH_{U',q}$'' --- the
      $g_2, \dots, g_m$ do not occur in the expansion displayed immediately above this
      sentence, $G' = \tau'(z_1+\alpha)^m + x_1\Omega_1 + y\Omega_2$. The phrase matches
      the parallel passage of the proof of Theorem 4.2 on printed p.\ 2516, where the
      expansion really has $g_2(z_1+\alpha)^{m-2} + \cdots + g_{m-1}(z_1+\alpha)$.
\item \textbf{p.\ 2520, first display:} ``$G' = g_h(z_1+\alpha)^{m-h}y + \cdots + g_m y +
      x_1\Omega_1 + y_1^2\Omega_2$'' --- the last term carries a subscript 1 on $y$, but
      this is the 1-point case with coordinates $x_1, y, z_1$; the parallel passage of the
      proof of Theorem 4.2 (printed p.\ 2516) prints $y\Omega_2$ there. Checked at
      1200\,dpi: the subscript 1 is really printed.
\item \textbf{p.\ 2521, last paragraph:} ``Since $2 \leqslant h \leqslant m-1$'' ---
      inherited from the parallel sentence on p.\ 2518; see the note in place. $h$ is
      defined there as the largest $i$ with $\tau_i \neq 0$ in (A), and (A) admits $i = m$.
      The displayed conclusion remains true for $h = m$ ($\sigma_D(q) = 0 < m-1$), so this
      is an unproved-but-harmless inherited bound, not a false statement.
\item \textbf{Ideal of the local ring: two spellings, both printed.} (c11) and (c12)
      print $I\OO_{U',q}$ (script O, no hat); the sentences that use principality inside
      the proof print $I\OH_{U',q}$ (with hat), and p.\ 2521's ``Suppose that
      $I\OO_{U',q}$ is principal'' returns to the unhatted form. Both kept exactly as
      printed.
\item \textbf{The surface in the last case: two spellings, both printed.} (c13) says
      ``the strict transform of $\ztil = 0$ in $U'$''; the ``final case'' paragraph says
      ``the strict transform of $z = 0$ in $U'$''. Both kept as printed.
\item \textbf{Re-use of the letters $a_1, b_1$:} in (c14) $x = x_1^{a_1}$,
      $z = x_1^{b_1}(z_1+\alpha)$; in (c15) $x = x_1^{a_1}z_1^{b_1}$,
      $z = x_1^{c_1}z_1^{d_1}$; in (c16) $z = x_1^{b_1}z_1$. These are the paper's own
      re-uses of the same letters across the three cases; kept as printed.
\item \textbf{Barred symbols:} $x_1^s z_1^{\tbar}$ (2-point case) and $b_1 t > \tbar$ ---
      the bar over $t$ is present in the image and is lost by the text layer (which prints
      a plain $t$). Transcribed with the bar.
\end{enumerate}

\subsection*{2.\ Reading record (which pages were read as images, and at what resolution)}
\begin{itemize}\itemsep0.2em
\item \textbf{pdftotext -layout} on the whole 18-page file (1036 lines) --- used ONLY to
      locate statements and to establish the paper's overall structure. The text layer is
      unusable for symbols here: it collapses $\OH$ to $\OO$, drops every overline/hat/tilde
      ($\Dbar \to D$, $\ztil \to z$, $\tbar \to t$), turns $I\OO_{U',q}$ and
      $I\OH_{U',q}$ into the same string, and renders $\fk$ and $\mathbb{Z}_+$ as ``k''
      and ``Z''.
\item \textbf{Read as images (read\_image), page by page:} PDF pp.\ 15, 16, 17, 18 at
      150\,dpi (whole pages) and at 300\,dpi (two half-page bands each); PDF p.\ 4 at
      300\,dpi (the crop containing item (4) and form (A), plus the line above it); PDF
      pp.\ 2 and 3 at 300\,dpi for the paper's own framing sentences that are quoted in
      this file (abstract and introduction on printed p.\ 2505, the opening of
      \S2 ``The corrections'' and the instruction ``Page 877: Definition 3.3 should be
      modified as follows.'' on printed p.\ 2506, together with the two errata items
      about $\omega'$ and the heading of the revised Definition 3.3); and
      these 600\,dpi or 1200\,dpi crops for symbol-level checks:
      (c11)/(c12)/(c13) on p.\ 2519; the $G$ expansion, the $G'$ expansion,
      $\ord(G'(0,0,z_1)) \leqslant m-1$ and $\sigma_{D'}(q) < m-1 = \sigma_D(p)$ on
      p.\ 2519; the line $G' = g_h(z_1+\alpha)^{m-h}y + \cdots + g_m y + x_1\Omega_1 +
      y_1^2\Omega_2$ and the sentence $\sigma_D(q) \leqslant (m-h+1)-1 < m-1 = \sigma_D(p)$
      on p.\ 2520; the sentences ``let $x_1^s z_1^{\tbar}$ be a local generator \dots\ We
      have that $a_1 t > s$ and $b_1 t > \tbar$'' and ``\dots\ we have that
      $\ord G'(0,y,0) \leqslant 1$'' on p.\ 2520; the whole of p.\ 2521's proof tail
      including the final $\square$; and the $\omega'$ paragraph opening on p.\ 2518.
      Every character of the collected unit was read from an image; no display formula was
      taken from the text layer.
\item \textbf{Crops kept locally:} \texttt{b21/pages/} (150\,dpi whole pages 15--18),
      \texttt{b21/img/} (300\,dpi bands \texttt{hd1*a/b}, 600\,dpi crops
      \texttt{z0,z1,z2,z3,z5,z6,z7}, 1200\,dpi crops \texttt{w1,w2,w3,w5,w6,w7}, the
      framing-sentence crops \texttt{intro-1\,\dots\,intro-4}, \texttt{def33-ann},
      \texttt{def33b\,\dots\,def33e}, and \texttt{p4-A}, \texttt{p4-pre},
      \texttt{p15-omega}, \texttt{p17-final}).
      Remote scratch: \texttt{/tmp/q2}, \texttt{/tmp/q2work} on \texttt{Yupeng-orx}.
      Kit: \texttt{pdflatex} (TeX Live 2025) compiled this file with exit~0 (6 pages).
\end{itemize}

\subsection*{3.\ What was deliberately NOT transcribed (and why)}
\begin{itemize}\itemsep0.2em
\item The errata's newly written \textbf{detailed proof of Theorem 4.2} (printed
      pp.\ 2515--2518 = PDF pp.\ 12--15). It is a complete proof supplied by this paper,
      but the \emph{statement} of Theorem 4.2 is not in this PDF: it lives in [3]
      (Ann.\ Inst.\ Fourier 63 (2013) 865--922, printed p.\ 898), and the errata only
      replaces the one-line proof of it. It is therefore not a self-contained
      theorem-plus-proof unit of this file. Only its tail (the last two lines, printed
      p.\ 2518) was read as an image, because CONTEXT B sits immediately below it, and its
      middle pages were read through the text layer for comparison with Theorem G's proof.
\item The revised \textbf{Lemma 3.6} and its proof (printed pp.\ 2507--2511) --- already
      collected under candidate P2 of this corpus; not re-transcribed here.
\item The revised \textbf{Lemma 3.7} (printed pp.\ 2511--2513): the errata changes its
      statement but does not re-prove it in this paper (its proof is in [3], with one
      inserted detailed sub-proof at ``Page 885, line 11''). Not a complete unit here.
\item The \textbf{line-level correction list} of the errata (printed pp.\ 2506, 2513,
      2515, 2521 telegraphed items), except the three items of CONTEXT C.
\item Nothing else: the whole of printed pp.\ 2518--2521 is transcribed above.
\end{itemize}

\end{document}
